\section{Related Work}

Our work intersects contamination detection, benchmark reliability, and learning dynamics analysis. We position against each area to show why existing approaches, while valuable, do not address contamination-performance quantification.

\subsection{Contamination Detection Methods}

\textbf{N-gram overlap detection.} The GPT-3 decontamination methodology \citep{brown2020language} established 13-gram matching as the standard for identifying verbatim overlap between training and evaluation data. EleutherAI's lm-evaluation-harness implements this approach for benchmark decontamination. \citet{yang2023rethinking} extended this analysis to RedPajama, finding 8--18\% overlap with HumanEval and demonstrating that simple paraphrasing bypasses n-gram detection. These methods detect contamination presence but do not quantify performance impact---a contaminated sample is flagged regardless of whether it affects model scores by 0.1\% or 10\%.

\textbf{Membership inference attacks.} MIA-based approaches \citep{carlini2021extracting,shi2024detecting} determine whether specific examples appeared in training data by analyzing model behavior signals: loss, perplexity, confidence distributions. The MIMIR benchmark \citep{duan2024mimir} standardizes these approaches for LLMs. Recent work on semantic membership inference \citep{mozaffari2024semantic} and keyword-based attacks \citep{antebi2025tagtab} improves detection accuracy. However, MIA methods answer ``was this seen?'' not ``how much did seeing it help?''

\textbf{Output distribution analysis.} \citet{dong2024generalization} proposed CDD (Contamination Detection via Distribution) to identify contaminated samples through output distribution peakedness, paired with TED (Test Data Deviation) for mitigation. \citet{choi2025contaminated} introduced Kernel Divergence Score for dataset-level contamination detection via fine-tuning sensitivity. These methods advance detection but remain focused on contamination presence rather than performance correlation.

\subsection{Benchmark Reliability}

\textbf{Contamination-resistant benchmarks.} Recognizing contamination concerns, recent work has developed frequently-updated benchmarks resistant to training data leakage. LiveBench \citep{white2024livebench} provides monthly-updated questions with objective ground-truth scoring. LessLeak-Bench \citep{zhou2025lessleak} covers 83 software engineering benchmarks with automated anti-leakage measures. These efforts sidestep contamination rather than quantifying its effects.

\textbf{Benchmark analysis.} \citet{sainz2023nlp} catalyzed the contamination discussion with ``NLP Evaluation in Trouble,'' documenting contamination levels across major benchmarks and arguing for community detection efforts. Their work establishes contamination as widespread but stops at detection---the correlation between contamination levels and score inflation remains unaddressed.

\subsection{Learning Dynamics and Memorization}

\textbf{Training dynamics analysis.} Research on learning dynamics \citep{swayamdipta2020dataset,toneva2019empirical} reveals that models learn different examples at different rates, with some samples memorized early while others require extended training. This work inspired our checkpoint-gradient approach---if memorization accumulates during training, contamination effects should correlate with training progress.

\textbf{Memorization studies.} \citet{carlini2023quantifying} demonstrated that LLMs memorize substantial training content, with extractable memorization scaling with model size and data repetition. \citet{biderman2023pythia} provided training dynamics for the Pythia model family with documented checkpoint availability---the transparency enabling our checkpoint-gradient methodology.

\subsection{Our Position}

Existing work establishes that contamination exists, can be detected, and that models memorize training content. What remains missing is the quantitative bridge: given X\% contamination, how much does the benchmark score inflate? We address this gap through checkpoint-gradient analysis with capability detrending, providing a methodology for contamination-inflation correlation measurement with preliminary validation results.
